\chapter{Modele o wartościach boolowskich i względna niesprzeczność CH}
\label{ch:boolean-valued-independence}
\chaptermark{Modele boolowskie i niesprzeczność}

\section{Dlaczego potrzebujemy innej konstrukcji?}

W Rozdziale~\ref{ch:forcing-continuum} wykazaliśmy, że z
\emph{istnienia przeliczalnego tranzytywnego modelu zbiorowego
ZFC} wynikają modele spełniające CH oraz modele spełniające
jej zaprzeczenie. Samo założenie
$\operatorname{Con}(\mathrm{ZFC})$ nie daje jednak takiego
modelu tranzytywnego. Poniżej uzasadnimy silniejszy wynik:
\[
\begin{aligned}
\operatorname{Con}(\mathrm{ZFC})
 &\Longrightarrow \operatorname{Con}(\mathrm{ZFC}+\mathrm{CH}),\\
\operatorname{Con}(\mathrm{ZFC})
 &\Longrightarrow \operatorname{Con}(\mathrm{ZFC}+\neg\mathrm{CH}).
\end{aligned}
\]
Argument będzie przebiegał przez modele o wartościach w
\emph{zupełnych algebrach Boole'a}. W przeciwieństwie do
poprzedniego rozdziału nie wymaga on, by model wyjściowy
był tranzytywny, dobrze ufundowany ani przeliczalny.

\begin{uwaga}[Metateoria i oznaczenia]
\label{rem:boolean35-metatheory}
Rozumujemy w zwykłej metateorii matematycznej opartej
na ZFC. Przez $M\models\mathrm{ZFC}$ rozumiemy dowolną
\emph{strukturę zbiorową} języka $\{\in\}$ spełniającą
wszystkie standardowe aksjomaty ZFC. Zbiór jej elementów
oznaczamy przez $|M|$; relację przynależności w $M$
przez $\in^M$. Nie zakładamy, że $\in^M$ pokrywa się z
zewnętrzną relacją $\in$. Jeśli $b\in|M|$ jest według $M$
algebrą Boole'a, przez $\mathbb B^M$ oznaczamy
\emph{zewnętrzny} zbiór jej elementów z działaniami
interpretowanymi przez $M$. Kresy w $\mathbb B$ są
zupełne \emph{w sensie $M$}; nie twierdzimy, że
istnieją kresy dowolnych zewnętrznych podzbiorów
$\mathbb B^M$.

Formuła $\varphi$ jest zawsze \emph{ustaloną standardową}
formułą pierwszego rzędu. Nie postulujemy pojedynczego
predykatu prawdy dla wszystkich wewnętrznych kodów
formuł w modelu $M$.
\end{uwaga}

\section{Od pojęcia wymuszania do algebry zupełnej}

\begin{definicja}[Algebra Boole'a i zupełność]
\label{def:boolean35-algebra}
Niech $\mathbb B$ będzie zbiorem z działaniami
$\wedge,\vee$, dopełnieniem $b\mapsto\neg b$
i elementami $0\ne1$, spełniającym aksjomaty
algebry Boole'a. Porządek określamy przez
$b\le c$ wtedy i tylko wtedy, gdy $b\wedge c=b$.
Algebra jest \emph{zupełna}, gdy każdy
podzbiór $A\subseteq\mathbb B$ ma kres górny
$\bigvee A$ i dolny $\bigwedge A$ w $\mathbb B$.
Implikację boolowską oznaczamy
$b\Rightarrow c:=\neg b\vee c$.
\end{definicja}

\begin{definicja}[Algebra zbiorów regularnie otwartych]
\label{def:boolean35-regular-open}
Niech $(X,\mathcal T)$ będzie przestrzenią
topologiczną. Zbiór otwarty $U\subseteq X$
nazywamy \emph{regularnie otwartym}, gdy
$U=\operatorname{int}(\overline U)$.
Rodzinę takich zbiorów oznaczamy
$\operatorname{RO}(X)$.
\end{definicja}

\begin{lemat}[Zupełność algebry zbiorów regularnie otwartych]
\label{lem:boolean35-ro-complete}
Dla każdej przestrzeni topologicznej
$(X,\mathcal T)$ zbiór $\operatorname{RO}(X)$
jest zupełną algebrą Boole'a z działaniami
\[
 U\wedge V=U\cap V,\qquad
 \neg U=\operatorname{int}(X\setminus\overline U),
 \qquad
 \bigvee_{i\in I}U_i=
 \operatorname{int}\overline{\bigcup_{i\in I}U_i}.
\]
Zbiór pusty i $X$ są odpowiednio $0$ i $1$.
\end{lemat}
\begin{proof}
Operator $r(A)=\operatorname{int}(\overline A)$
jest idempotentny na zbiorach otwartych:
dla otwartego $U$ zachodzi
$\overline{r(U)}=\overline U$, gdyż
$U\subseteq r(U)\subseteq\overline U$.
Stąd $r(r(U))=r(U)$.
Jeżeli $U,V$ są regularnie otwarte,
to $U\cap V$ jest otwarty i
\[
 U\cap V\subseteq r(U\cap V)
 \subseteq r(U)\cap r(V)=U\cap V.
\]
Wskazany wzór na dopełnienie daje
zbiór regularnie otwarty rozłączny z $U$;
suma boolowska obu jest równa $X$,
ponieważ ich suma jest gęsta w $X$.
Wzór na dowolny kres górny jest
najmniejszym regularnie otwartym
zbiorem zawierającym wszystkie $U_i$.
Kres dolny dowolnej rodziny definiujemy
z praw de Morgana; dla pustej rodziny
są to odpowiednio $0$ i $1$.
Pozostałe tożsamości boolowskie
wynikają z powyższych działań,
w szczególności z rozdzielności
przecięcia względem regularnego
domknięcia sumy zbiorów otwartych.
\end{proof}

\begin{definicja}[Uzupełnienie boolowskie wymuszania]
\label{def:boolean35-completion}
Niech $(\mathbb P,\leq)$ będzie niepustym
pojęciem wymuszania. Na $\mathbb P$
rozpatrujemy topologię, której bazą są
stożki $\mathbb P_{\le p}=\{q:q\le p\}$.
Algebrę
$\mathbb B_{\mathbb P}=\operatorname{RO}(\mathbb P)$
nazywamy jego \emph{uzupełnieniem boolowskim}.
Warunkowi $p$ odpowiada niezerowy element
$i(p)=\operatorname{int}\overline{\mathbb P_{\le p}}$.
Przy porównywaniu wymuszań korzystamy z
dodatnich elementów $\mathbb B_{\mathbb P}$
i ich porządku; równoważne warunki
preorderu mogą dawać ten sam element $i(p)$.
\end{definicja}

\begin{lemat}[Gęstość obrazów warunków]
\label{lem:boolean35-dense-completion}
Dla niepustego pojęcia wymuszania
$(\mathbb P,\leq)$ oraz algebry
$\mathbb B_{\mathbb P}$ z
Definicji~\ref{def:boolean35-completion}
zbiór $\{i(p):p\in\mathbb P\}$
jest gęsty w $\mathbb B_{\mathbb P}^+
=\mathbb B_{\mathbb P}\setminus\{0\}$:
dla każdego $b>0$ istnieje
$p\in\mathbb P$ z $0<i(p)\le b$.
Jeśli $p,q$ mają wspólne wzmocnienie,
to $i(p)\wedge i(q)>0$; jeżeli są niezgodne,
to $i(p)\wedge i(q)=0$.
\end{lemat}
\begin{proof}
Niech $b$ będzie niepustym regularnie
otwartym zbiorem warunków. Wybieramy
$p\in b$. Ponieważ $b$ jest otwarty,
istnieje $q\le p$ takie, że
$\mathbb P_{\le q}\subseteq b$.
Z regularnej otwartości $b$ wynika
$i(q)\subseteq b$. Stożek nie jest pusty,
więc $i(q)>0$. Jeżeli istnieje
$r\le p,q$, to $r$ należy do przecięcia
stożków i jego regularne otwarte
otoczenie świadczy o niezerowości
$i(p)\cap i(q)$. Gdy wspólnego
wzmocnienia nie ma, stożki są rozłączne,
a przecięcie ich regularnych otwartych
domknięć jest puste: gdyby było
niepuste, zawierałoby stożek i na nim
oba pierwotne stożki byłyby gęste,
co dostarczyłoby wspólnego wzmocnienia.
\end{proof}

\section{Nazwy boolowskie i wartości formuł}

\begin{definicja}[Hierarchia nazw boolowskich]
\label{def:boolean35-names}
Niech $\mathbb B$ będzie niezerową
zupełną algebrą Boole'a.
Rekurencyjnie po liczbach porządkowych
definiujemy
\[
 V_0^{\mathbb B}=\varnothing,\quad
 V_{\alpha+1}^{\mathbb B}
 =\{f:\operatorname{dom}f\subseteq
   V_\alpha^{\mathbb B},\
   f:\operatorname{dom}f\to\mathbb B\},
 \quad
 V_\lambda^{\mathbb B}
 =\bigcup_{\alpha<\lambda}V_\alpha^{\mathbb B}.
\]
Elementy klasy
$V^{\mathbb B}=\bigcup_{\alpha\in\operatorname{Ord}}
V_\alpha^{\mathbb B}$ nazywamy
\emph{nazwami boolowskimi}.
Dla $\tau\in V^{\mathbb B}$ zapis
$\tau(\sigma)$ oznacza współczynnik
$\sigma\in\operatorname{dom}\tau$.
Dla $x$ z modelu wyjściowego
nazwa kanoniczna $\check x$ ma dziedzinę
$\{\check y:y\in x\}$ i przypisuje każdemu
$\check y$ wartość $1$.
\end{definicja}

\begin{definicja}[Wartości boolowskie formuł]
\label{def:boolean35-values}
Dla nazw $\sigma,\tau$ definiujemy
jednoczesną rekurencję atomową
\[
\begin{aligned}
\llbracket\sigma\in\tau\rrbracket
 &=\bigvee_{\pi\in\operatorname{dom}\tau}
       \bigl(\tau(\pi)\wedge
             \llbracket\sigma=\pi\rrbracket\bigr),\\
\llbracket\sigma=\tau\rrbracket
 &=\bigwedge_{\pi\in\operatorname{dom}\sigma}
       \bigl(\sigma(\pi)\Rightarrow
             \llbracket\pi\in\tau\rrbracket\bigr)
   \wedge
   \bigwedge_{\pi\in\operatorname{dom}\tau}
       \bigl(\tau(\pi)\Rightarrow
             \llbracket\pi\in\sigma\rrbracket\bigr).
\end{aligned}
\]
Kolejne klauzule dla ustalonych formuł są:
\[
\begin{aligned}
\llbracket\neg\varphi\rrbracket
   &=\neg\llbracket\varphi\rrbracket,\qquad
\llbracket\varphi\land\psi\rrbracket
   =\llbracket\varphi\rrbracket
       \wedge\llbracket\psi\rrbracket,\\
\llbracket\exists x\,\varphi(x)\rrbracket
   &=\bigvee\{\llbracket\varphi(\tau)\rrbracket:
               \tau\in V^{\mathbb B}\}.
\end{aligned}
\]
Pozostałe spójniki i kwantyfikator
ogólny definiujemy zwyczajowo.
\end{definicja}

\begin{lemat}[Poprawność definicji wartości]
\label{lem:boolean35-well-defined}
Niech $\mathbb B$ będzie niezerową
zupełną algebrą Boole'a.
Dla każdej ustalonej formuły
$\varphi(x_1,\ldots,x_n)$ i wszystkich
nazw $\tau_1,\ldots,\tau_n$
wartość $\llbracket
\varphi(\tau_1,\ldots,\tau_n)\rrbracket$
jest jednoznacznie określonym
elementem $\mathbb B$. Operacja ta
jest definiowalna w ZFC przy
ustalonej formule $\varphi$.
\end{lemat}
\begin{proof}
Klauzule atomowe są poprawnie
ufundowane przez rekurencję
po \emph{sumie naturalnej rang}
dwóch nazw, podobnie jak w
Uwadze~\ref{rem:forcing33-ranks}.
Przy każdej rekurencyjnej zależności
jedna z nazw jest zastępowana nazwą
ze zbioru wartości współrzędnych
pierwszej nazwy, o ściśle mniejszej
randze. Kresy w obu klauzulach
dotyczą zbiorów, a ich istnienie
zapewnia zupełność $\mathbb B$.

Pozostaje uzasadnić kwantyfikator
egzystencjalny. Choć
$V^{\mathbb B}$ jest klasą właściwą,
\emph{zbiór różnych wartości}
\[
 S_\varphi=\{b\in\mathbb B:
      \exists\tau\ (\tau\text{ jest nazwą}
        \land b=\llbracket\varphi(\tau)\rrbracket)\}
\]
jest podzbiorem \emph{zbioru}
$\mathbb B$. Powstaje przez
wyróżnianie, gdy klauzula dla
$\varphi$ jest już ustalona.
Zupełność $\mathbb B$ daje
$\bigvee S_\varphi$.
Indukcja po budowie każdej
ustalonej formuły kończy dowód.
Nie twierdzimy, że zbiór
wszystkich nazw jest elementem
uniwersum.
\end{proof}

\begin{lemat}[Tożsamości równości i podstawiania]
\label{lem:boolean35-equality}
Niech $\mathbb B$ będzie niezerową
zupełną algebrą Boole'a i
$\sigma,\tau,\rho$ nazwami.
Wtedy
\[
\begin{aligned}
\llbracket\sigma=\sigma\rrbracket&=1,\\
\llbracket\sigma=\tau\rrbracket
 &=\llbracket\tau=\sigma\rrbracket,\\
\llbracket\sigma=\tau\rrbracket\wedge
 \llbracket\tau=\rho\rrbracket
 &\le\llbracket\sigma=\rho\rrbracket.
\end{aligned}
\]
Dla każdej ustalonej formuły $\varphi(x)$
z dodatkowymi parametrami nazwami mamy
\[
 \llbracket\sigma=\tau\rrbracket\wedge
 \llbracket\varphi(\sigma)\rrbracket
 \le\llbracket\varphi(\tau)\rrbracket.
\]
\end{lemat}
\begin{proof}
Refleksyjność, symetrię i przechodniość
weryfikujemy równoczesną indukcją
po rangach nazw, korzystając z
rozdzielności zupełnej algebry Boole'a:
$b\wedge\bigvee_i c_i
=\bigvee_i(b\wedge c_i)$.
Dla przechodniości najpierw
przechodzimy od współczynnika
$\sigma(\pi)$ do przynależności
$\pi\in\tau$, a następnie do
$\pi\in\rho$; nierówności
$\llbracket\sigma=\tau\rrbracket
 \wedge\llbracket\tau=\rho\rrbracket
 \le\llbracket\sigma=\rho\rrbracket$
wynikają z dwóch klauzul
ograniczających współczynniki
równości. Analogicznie dla
$\llbracket\sigma\in\tau\rrbracket$.

Podstawianie dowodzimy indukcją
po budowie $\varphi$. W przypadku
atomowym stosujemy poprzednie
nierówności i definicję wartości
przynależności. Koniunkcja
zachowuje nierówność przez
monotoniczność $\wedge$, negacja
przez symetrię równości i
odwrócenie porządku przy
dopełnieniu. Dla egzystencji
stosujemy rozdzielność
$\wedge$ względem dowolnych
kresów górnych.
\end{proof}

\section{Mieszanie nazw i lemat o świadkach}

\begin{definicja}[Antyłańcuch w algebrze Boole'a]
\label{def:boolean35-antichain}
Niech $\mathbb B$ będzie algebrą Boole'a.
Zbiór $A\subseteq\mathbb B\setminus\{0\}$
jest \emph{antyłańcuchem boolowskim},
jeżeli $a\wedge a'=0$ dla wszystkich
różnych $a,a'\in A$. Jest maksymalny
poniżej $b\in\mathbb B$, gdy
$a\le b$ dla każdego $a\in A$
oraz $\bigvee A=b$.
\end{definicja}

\begin{lemat}[Mieszanie nazw]
\label{lem:boolean35-mixing}
Niech $\mathbb B$ będzie zupełną
algebrą Boole'a, $A\subseteq\mathbb B^+$
antyłańcuchem i
$(\tau_a)_{a\in A}$ rodziną nazw
boolowskich \emph{indeksowaną zbiorem} $A$.
Wtedy istnieje nazwa $\tau$ taka, że
\[
 a\le\llbracket\tau=\tau_a\rrbracket
 \qquad(a\in A).
\]
\end{lemat}
\begin{proof}
Niech $D=\bigcup_{a\in A}
\operatorname{dom}\tau_a$. Zastępowanie
i suma zapewniają, że $D$ jest zbiorem
nazw o wspólnie ograniczonych rangach.
Definiujemy na $D$ funkcję
\[
 \tau(\pi)=
  \bigvee\{a\wedge\tau_a(\pi):
       a\in A,\ \pi\in\operatorname{dom}\tau_a\}.
\]
Jest to nazwa, ponieważ wszystkie
jej współczynniki należą do $\mathbb B$.
Dla ustalonego $a\in A$ i $\pi\in D$
rozdzielność i rozłączność
antyłańcucha dają
$a\wedge\tau(\pi)=
 a\wedge\tau_a(\pi)$, gdy
$\pi\in\operatorname{dom}\tau_a$,
oraz $a\wedge\tau(\pi)=0$ w przeciwnym
przypadku. Wstawiając te równości
do atomowych klauzul
Definicji~\ref{def:boolean35-values}
i korzystając z refleksyjności
równości z Lematu~\ref{lem:boolean35-equality},
otrzymujemy
$a\le\llbracket\tau=\tau_a\rrbracket$.
\end{proof}

\begin{twierdzenie}[Lemat o pełności wartości]
\label{thm:boolean35-fullness}
Niech $\mathbb B$ będzie niezerową
zupełną algebrą Boole'a.
Dla każdej ustalonej formuły
$\varphi(x,\tau_1,\ldots,\tau_n)$
oraz nazw $\tau_1,\ldots,\tau_n$
istnieje nazwa $\rho$ spełniająca
\[
 \llbracket\exists x\,\varphi(x,\vec\tau)\rrbracket
   =\llbracket\varphi(\rho,\vec\tau)\rrbracket.
\]
\end{twierdzenie}
\begin{proof}
Niech $u=\bigvee S_\varphi$, gdzie
$S_\varphi\subseteq\mathbb B$
jest zbiorem wartości
z Lematu~\ref{lem:boolean35-well-defined}.
Jeżeli $u=0$, dowolna nazwa
jest świadkiem, ponieważ jej wartość
nie przekracza $u$. Załóżmy $u>0$.
Zdefiniujmy
\[
 D=\{a\in\mathbb B^+:
    a\le u,\ \exists b\in S_\varphi\ (a\le b)\}.
\]
Zbiór $D$ jest gęsty poniżej $u$:
jeżeli $0<c\le u$, to z
$c\wedge\bigvee S_\varphi=c>0$
i rozdzielności istnieje
$b\in S_\varphi$ z $c\wedge b>0$;
wówczas $c\wedge b\in D$.
Wybierzmy maksymalny antyłańcuch
$A\subseteq D$, którego kres górny
jest równy $u$.

Dla każdego $a\in A$ istnieje
nazwa $\rho_a$ taka, że
$a\le\llbracket\varphi(\rho_a,\vec\tau)
\rrbracket$. Nie stosujemy wyboru
z klasy właściwej: dla każdego $a$
wybieramy \emph{najmniejszą rangę}
nazwy spełniającej nierówność.
Istnienie minimum wynika ze
znalezienia jednej nazwy i
minimalizacji w zbiorze liczb
porządkowych poniżej jej rangi.
Zastępowanie na zbiorze $A$
ogranicza wszystkie minimalne
rangi wspólną liczbą porządkową
$\gamma$. Teraz wszystkie
kandydatury leżą w jednym
\emph{zbiorze} $V_{\gamma+1}^{\mathbb B}$;
zastosowanie aksjomatu wyboru
daje rodzinę $(\rho_a)_{a\in A}$.

Lemat~\ref{lem:boolean35-mixing}
daje nazwę $\rho$ równą $\rho_a$
z wartością co najmniej $a$.
Z Lematu~\ref{lem:boolean35-equality}
i podstawiania wynika
\[
 a\le\llbracket\varphi(\rho,\vec\tau)\rrbracket
 \quad(a\in A).
\]
Ponieważ $\bigvee A=u$,
mamy $u\le\llbracket\varphi(\rho,\vec\tau)
\rrbracket$. Odwrotna nierówność
wynika z definicji $u$ jako kresu
wszystkich wartości formuły.
\end{proof}

\section{Aksjomaty ZFC w uniwersum boolowskim}

\begin{twierdzenie}[Boolowska poprawność ZFC]
\label{thm:boolean35-zfc}
Niech $\mathbb B$ będzie
niezerową zupełną algebrą Boole'a.
Wtedy dla każdego \emph{ustalonego}
aksjomatu ZFC oraz każdego
ustalonego wystąpienia schematów
wyróżniania i zastępowania,
oznaczonego przez $\theta$,
zachodzi
$\llbracket\theta\rrbracket=1$.
Dotyczy to także wszystkich
aksjomatów logicznych i
reguł poprawnego wnioskowania.
\end{twierdzenie}
\begin{proof}
Poniżej wskazujemy konstrukcje nazw
i uzasadnienie wartości $1$ dla
każdej grupy aksjomatów. W całym
dowodzie korzystamy z
Lematu~\ref{lem:boolean35-equality},
w szczególności z podstawiania.

\emph{Ekstensjonalność.}
Klauzula równości stwierdza, że
każdy element nazwy $\sigma$,
z odpowiednim współczynnikiem,
należy do $\tau$ i odwrotnie.
Użycie obu implikacji boolowskich
oraz podstawiania daje
\[
 \llbracket\forall x\,
 (x\in\sigma\leftrightarrow x\in\tau)
 \rrbracket
 \le\llbracket\sigma=\tau\rrbracket.
\]
W drugą stronę stosujemy podstawianie.
Zatem ekstensjonalność ma wartość $1$.

\emph{Pusty zbiór, para i nieskończoność.}
Pusta funkcja jest nazwą zbioru
pustego. Nazwa o współczynnikach
$1$ przy $\sigma,\tau$ reprezentuje
$\{\sigma,\tau\}$, z uwzględnieniem
możliwego zlania współrzędnych,
gdy nazwy są identyczne.
Nazwa $\check\omega$ reprezentuje
wewnętrznie zbiór induktywny;
definicja $\check x$ zachowuje
przynależność i podstawowe
konstrukcje zbiorowe z wartością $1$.

\emph{Aksjomat sumy.}
Dla nazwy $\sigma$ ustalamy
$D=\bigcup_{\pi\in\operatorname{dom}\sigma}
\operatorname{dom}\pi$. To jest zbiór.
Tworzymy na $D$ nazwę
\[
 U_\sigma(\rho)=
 \bigvee_{\pi\in\operatorname{dom}\sigma}
   \bigl(\sigma(\pi)\wedge
     \pi(\rho)\bigr),
\]
przyjmując $\pi(\rho)=0$ dla
$\rho\notin\operatorname{dom}\pi$.
Klauzule atomowe, rozdzielność
kresów i podstawianie pokazują,
że dla każdej nazwy $\eta$
\[
 \llbracket\eta\in U_\sigma\rrbracket
 =\llbracket
    \exists x\in\sigma\ (\eta\in x)\rrbracket.
\]
To jest aksjomat sumy z wartością $1$.

\emph{Schemat wyróżniania.}
Niech $\varphi(x,\vec\tau)$ będzie
ustaloną formułą, a $\sigma$ nazwą.
Na zbiorze $\operatorname{dom}\sigma$
definiujemy nazwę
\[
 S(\rho)=\sigma(\rho)\wedge
      \llbracket\varphi(\rho,\vec\tau)\rrbracket.
\]
Podstawianie oraz rozdzielność
dają dla każdej nazwy $\eta$
\[
 \llbracket\eta\in S\rrbracket
 =\llbracket\eta\in\sigma\land
                  \varphi(\eta,\vec\tau)\rrbracket.
\]
To potwierdza każdy ustalony
przypadek schematu wyróżniania.

\emph{Aksjomat zbioru potęgowego.}
Niech $D=\operatorname{dom}\sigma$.
W otaczającym modelu wyjściowym
rozpatrujemy \emph{zbiór} wszystkich
funkcji $f:D\to\mathbb B$ takich,
że $f(\rho)\le\sigma(\rho)$
dla każdej $\rho\in D$.
Każda taka funkcja jest nazwą.
Tworzymy nazwę $\Pi_\sigma$,
która każdej z tych funkcji
przypisuje wartość $1$.
Jeżeli $f$ spełnia wskazane
ograniczenia, podstawianie daje
$\llbracket f\subseteq\sigma\rrbracket=1$.

Odwrotnie, niech $\tau$ będzie
dowolną nazwą i niech
$b=\llbracket\tau\subseteq\sigma\rrbracket$.
Definiujemy funkcję na $D$
przez
\[
 f_\tau(\rho)=
   \sigma(\rho)\wedge
   \llbracket\rho\in\tau\rrbracket.
\]
Z aksjomatu ekstensjonalności
w postaci boolowskiej wynika
$b\le\llbracket\tau=f_\tau\rrbracket$:
pod wartością $b$ każdy element
$\tau$ jest elementem $\sigma$,
a jego przynależność można więc
odczytać przez współczynniki
na ustalonej dziedzinie $D$.
Wartość $b$ nie przekracza zatem
$\llbracket\tau\in\Pi_\sigma\rrbracket$.
Obie strony definicji zbioru
potęgowego mają wartość $1$.

\emph{Schemat zastępowania.}
Niech $\varphi(x,y,\vec\tau)$
będzie ustaloną formułą, a
$\sigma$ nazwą. Jeżeli
\[
 b=\llbracket
  \forall x\in\sigma\ \exists!y\
            \varphi(x,y,\vec\tau)\rrbracket,
\]
to dla każdego
$\pi\in\operatorname{dom}\sigma$
Lemat~\ref{thm:boolean35-fullness}
dostarcza nazwy $\rho_\pi$
uzyskującej wartość
$\llbracket\exists y\,
\varphi(\pi,y,\vec\tau)\rrbracket$.
Rodzinę świadków dla \emph{zbioru}
$\operatorname{dom}\sigma$ wybieramy
za pomocą minimalnych rang,
zastępowania i wyboru,
tak samo jak w dowodzie
Lematu~\ref{thm:boolean35-fullness}.
Z nazw $\rho_\pi$ tworzymy
nazwę obrazu o współczynnikach
$\sigma(\pi)$, łącząc powtarzające
się nazwy przez kres górny.
Pod wartością $b$ każde
$x\in\sigma$ ma jedyny obraz
spełniający $\varphi$, a wskazana
nazwa zawiera dokładnie te obrazy.
Dla każdej ustalonej formuły
jest to wymagany przypadek
schematu zastępowania.

\emph{Aksjomat ufundowania.}
Ustalmy nazwę $\sigma$.
Jeżeli $b=\llbracket
\sigma\ne\varnothing\rrbracket>0$,
to pod każdym niezerowym
$c\le b$ istnieje nazwa $\pi$
i niezerowe $d\le c$ takie,
że $d\le\llbracket\pi\in\sigma\rrbracket$.
Wybieramy na odpowiednim
niezerowym fragmencie
\emph{najmniejszą możliwą rangę}
takiego reprezentanta.
Ta minimalizacja jest poprawna,
bo każdą przynależność do $\sigma$
można reprezentować nazwą z
\emph{zbioru} $\operatorname{dom}\sigma$;
zbiór możliwych rang jest więc
ograniczony. Jeśli na niezerowym
fragmencie istniałby jednocześnie
element $\rho\in\pi\cap\sigma$,
atomowa klauzula przynależności
do $\pi$ dawałaby, po ewentualnym
zastąpieniu $\rho$ równą jej nazwą
ze współrzędnych $\pi$, nazwę
o \emph{mniejszej randze}
reprezentującą element $\sigma$
na tym fragmencie. To przeczy
wyborowi najmniejszej rangi.
Wybór maksymalnego antyłańcucha
takich fragmentów i
Lemat~\ref{lem:boolean35-mixing}
sklejają lokalne świadectwa
w jedną nazwę elementu
$\pi\in\sigma$ rozłącznego
z $\sigma$ z wartością co
najmniej $b$. Stąd
\[
 \llbracket\forall x\,
   (x\ne\varnothing\Longrightarrow
    \exists y\in x\ (x\cap y=\varnothing))
 \rrbracket=1.
\]

\emph{Aksjomat wyboru.}
Dla nazwy $\sigma$ zbiór
$D=\operatorname{dom}\sigma$
ma w modelu wyjściowym dobre
uporządkowanie, a więc surjekcję
$e:\kappa\to D$ dla pewnej
liczby porządkowej $\kappa$
(gdy $D$ jest pusty, wybieramy
$\kappa=0$). Nazwa funkcji
$\xi\mapsto e(\xi)$ ma z wartością
$1$ obraz zawierający wszystkie
elementy $\sigma$. Wartości
$\llbracket e(\xi)\in\sigma\rrbracket$
wyznaczają nazwę podzbioru
$\kappa$ indeksów faktycznie
należących do $\sigma$.
Ograniczenie tej funkcji do
owej nazwanej dziedziny jest
surjekcją na $\sigma$ z
wartością $1$. Ponieważ
dziedzinę można dobrze
uporządkować, każdy element
$\sigma$ otrzymuje najmniejszy
indeks swojego przeciwobrazu.
Wcześniej zweryfikowane
wyróżnianie i zastępowanie
pozwalają przenieść to dobre
uporządkowanie na $\sigma$.
Aksjomat wyboru ma więc
wartość $1$.

Pozostają logiczne reguły
rachunku pierwszego rzędu.
Wszystkie aksjomaty logiczne
mają wartość $1$ na mocy
tożsamości algebry Boole'a,
a modus ponens zachowuje
wartość $1$. Reguła uogólniania
zachowuje wartość $1$, ponieważ
kres dolny rodziny wartości
równych $1$ jest równy $1$.
Jest to poprawność boolowska
całego rachunku.
\end{proof}
